% Results (main text). Numbers belong in the figures, tables and captions; the text states what they show and why it
% matters. The numbers behind every statement are in the Supplementary Results and Tables (supplement.tex).
\section{Results}
\label{sec:results}

\subsection{\sone{} files are the smallest on every matrix}

\sone{} produced the smallest file on every one of the \nm{nMatrices} matrices (Fig.~\ref{fig:wow}a,
Table~\ref{tab:main}). The typical gzip-compressed AnnData file, the most common way single-cell data are shared, was
\nm{ratio-h5ad_gzip} times larger, and even the smallest alternative on each matrix was \nm{ratioBestOther} times larger.
The gain does not come from a stronger general-purpose compressor: zstd at its strongest level on the raw arrays, and a
cloud-ready Zarr layout with the same block size as \sone, both left files about twice as large. It comes from
describing the counts. Measured against the floor of the best count model, three quarters of a gzip-compressed AnnData
file is redundant, and at least \nm{redund-other-min}\% of every other format; for \sone{} the median is
\nm{redund-s1}\% (Table~\ref{tab:main}).

% Table 2 (main text): every format on the same matrices, from generated/tab_main.tsv (report_limits.redundancy_section).
\begin{table*}[tbp]
\centering
\caption{Every format on the same \nm{sumNsize} matrices}
\label{tab:main}
\footnotesize\setlength{\tabcolsep}{5pt}
\pgfplotstabletypeset[
  columns={format,bits,rel,cr,read,fastest,read1,fastest1,redund},
  every head row/.style={before row={\toprule
    & \multicolumn{3}{c}{Size} & \multicolumn{2}{c}{Four cores} & \multicolumn{2}{c}{One core} & \\
    \cmidrule(lr){2-4}\cmidrule(lr){5-6}\cmidrule(lr){7-8}}, after row=\midrule},
  every last row/.style={after row=\bottomrule},
  columns/format/.style={string type, column name=Format, column type=l},
  columns/bits/.style={column name={Bits per non-zero}, fixed, fixed zerofill, precision=1},
  columns/rel/.style={column name={Size $\div$ \sone}, fixed, fixed zerofill, precision=1, postproc cell content/.append style={/pgfplots/table/@cell content/.add={}{$\times$}}},
  columns/cr/.style={column name={Compression}, fixed, fixed zerofill, precision=1, postproc cell content/.append style={/pgfplots/table/@cell content/.add={}{$\times$}}},
  columns/read/.style={column name={Read rate}, fixed, precision=0},
  columns/fastest/.style={string type, column name=Fastest, column type=r},
  columns/read1/.style={column name={Read rate}, fixed, precision=0},
  columns/fastest1/.style={string type, column name=Fastest, column type=r},
  columns/redund/.style={column name={Redundancy}, fixed, precision=0, postproc cell content/.append style={/pgfplots/table/@cell content/.add={}{\%}}},
]{\gen{tab_main.tsv}}
\par\vspace{3pt}
\begin{minipage}{\linewidth}\scriptsize
Bits per non-zero: median over matrices, for the whole file with gene identifiers and barcodes. Size $\div$ \sone:
geometric mean of the file's size relative to the \sone{} file of the same matrix. Compression: median ratio of the size
of uncompressed CSC arrays with 32-bit indices and values to the file's size. Read rate: median cold read rate into an
in-memory sparse matrix, million non-zeros per second, on four cores and on one core. Fastest: matrices on which the
format read fastest (one core: the \nm{dec-paired1-n} matrices measured on one core for every format). Redundancy: the
share of the file that a code as short as the best count model we fitted would not need, $1-L/S$, where $S$ is the file's
size and $L$ the cross-validated code length of that model on the same matrix (Section~\ref{sec:floor}); median over
matrices. R formats hold the \nm{sumNr} matrices below $2^{31}$ non-zeros and loom the \nm{sumNloom} it can store
densely.
\end{minipage}
\end{table*}


\subsection{Smaller files that are also faster to read}

A smaller file usually costs more time to decompress. \sone{} does not pay that price: read from a cold file into a sparse
matrix, it was the fastest format on \nm{dec-s1-fastest-count} of \nm{dec-paired-n} matrices and the only one that was
both small and fast (Fig.~\ref{fig:wow}b). Part of the reason is that its decoder spreads the work over the available
cores. Most formats do not decode in parallel at all; among the few that do, \sone{} gained the most from each added core
(Fig.~\ref{fig:wow}d).


\subsection{\sone{} reads at the speed of the network}
\label{sec:network}

On four cores \sone{} keeps up with links of about 1\,Gbit/s, and with faster links on more cores. Because each entry
takes so few bits, the same link then carries about twelve times as many entries as \sone{} as it does as uncompressed
arrays, and more than any other format over every common link speed (Fig.~\ref{fig:net}a).

Real training loaders, reading 512-cell minibatches into PyTorch from an object store whose bandwidth we controlled,
showed the same picture (Fig.~\ref{fig:net}b). The \sone{} loader filled most of a 1\,Gbit/s link and delivered thirty to
a hundred times more cells per second than AnnData's Zarr loader or TileDB-SOMA's loader for machine learning, and its
first minibatch arrived several times sooner than from any other loader that reads a file in parts (Fig.~\ref{fig:wow}c).
Latency matters as much as bandwidth: every request waits a full round trip, and cloud storage adds tens of milliseconds to
each. Because the \sone{} index locates every block before any is read, the reader sends all of a minibatch's requests at
once and keeps the next minibatches' requests in flight. Our measurements showed that eight concurrent requests left a
high-latency link mostly idle; with more in flight the rate recovered (Fig.~\ref{fig:net}c), and \texttt{singlet} 1.3.3 now
uses these settings by default.

Finally, we read the files from a real cloud object store, Cloudflare R2, over the public internet
(Fig.~\ref{fig:net}d). The first minibatch of a 580,000-cell lung atlas arrived in under a second, while downloading the
AnnData file first took well over a minute, and \sone{} then streamed cells tens of times faster than either Zarr layout.
Because it moves fewer bytes in fewer requests, an epoch of random minibatches would cost about twenty times less in
transfer and request fees than with AnnData Zarr, at the prices of a provider that charges for both. \sone{} needs no
extra DataLoader worker processes to reach these rates, since its decoder already uses every core.

% Network figure: what "at the speed of the network" means, and which formats and readers get there. Every value is read from
% generated/fignet_*.tsv (written by `python -m s1bench report`, network_section).
\begin{figure*}[p]
\centering
% ---------------- (a) delivered rate against link bandwidth ----------------
\begin{tikzpicture}
\begin{axis}[
  title={a\quad Delivered rate against link bandwidth},
  width=0.41\linewidth, height=6.8cm, xmode=log, ymode=log, xmin=0.01, xmax=100, ymin=0.3, ymax=1500,
  xlabel={link bandwidth $B$ (Gbit/s)}, ylabel={non-zeros delivered, millions per second},
  xtick={0.01,0.1,1,10,100}, xticklabels={0.01,0.1,1,10,100}, axis x line*=bottom, axis y line*=left,
  every axis plot/.append style={no marks},
  clip=false,
]
\pgfplotsinvokeforeach{0.1,1,10}{\draw[gen!35, line width=0.4pt] (axis cs:#1,0.3) -- (axis cs:#1,1500);}
% every format, thin grey; then the ones named in the text
\pgfplotsinvokeforeach{mtx_gz,tenx_h5,h5ad_none,h5ad_lzf,loom,soma,parquet_coo,rds_gz}{
  \addplot[gen!45, line width=0.5pt] table[x=gbit, y=#1] {\gen{fignet_curves.tsv}};}
\addplot[ink2, thick] table[x=gbit, y=csc_i32] {\gen{fignet_curves.tsv}} node[pos=1, anchor=west, font=\sffamily\tiny] {CSC, int32};
\addplot[dl, thick] table[x=gbit, y=zarr] {\gen{fignet_curves.tsv}} node[pos=1, anchor=west, font=\sffamily\tiny, text=dl] {Zarr};
\addplot[tens, thick, dashed] table[x=gbit, y=csc_zstd19] {\gen{fignet_curves.tsv}} node[pos=1, anchor=west, font=\sffamily\tiny, text=tens] {zstd-19};
\addplot[dec, thick] table[x=gbit, y=h5ad_gzip] {\gen{fignet_curves.tsv}} node[pos=1, anchor=west, font=\sffamily\tiny, text=dec!80!black] {h5ad gzip};
\addplot[red!65!black, thick, dashed] table[x=gbit, y=rds_xz] {\gen{fignet_curves.tsv}} node[pos=1, anchor=west, font=\sffamily\tiny, text=red!65!black] {rds xz};
\addplot[dom, thick] table[x=gbit, y=bpcells] {\gen{fignet_curves.tsv}} node[pos=1, anchor=west, font=\sffamily\tiny, text=dom] {BPCells};
\addplot[sone, densely dotted] table[x=gbit, y=s1_net] {\gen{fignet_curves.tsv}}
  node[pos=0.22, sloped, above, font=\sffamily\tiny, text=sone] {network-bound: $B/b$};
\addplot[sone, densely dotted] table[x=gbit, y=s1_dec] {\gen{fignet_curves.tsv}}
  node[pos=0.02, anchor=south west, font=\sffamily\tiny, text=sone] {decoder-bound: $D_p$};
\addplot[sone, line width=1.6pt] table[x=gbit, y=s1] {\gen{fignet_curves.tsv}} node[pos=1, anchor=west, font=\sffamily\tiny\bfseries, text=sone] {.s1};
\pgfplotsinvokeforeach{0.1,1,10}{\node[font=\sffamily\tiny, text=gen, anchor=south west, inner sep=1pt] at (axis cs:#1,0.32) {#1};}
\end{axis}
\end{tikzpicture}\hfill
% ---------------- (b) the link speed each decoder keeps up with ----------------
\begin{tikzpicture}
\begin{axis}[
  title={b\quad Which links each format's decoder keeps up with},
  width=0.47\linewidth, height=6.8cm, xmode=log, ymode=log, xmin=3, xmax=100, ymin=0.03, ymax=60,
  xlabel={stored bits per non-zero $b$}, ylabel={stored bits decoded per second, $D_p\,b$ (Gbit/s)},
  xtick={5,10,20,50,100}, xticklabels={5,10,20,50,100}, ytick={0.1,1,10}, yticklabels={0.1,1,10},
  axis x line*=bottom, axis y line*=left, clip mode=individual,
  lbl/.style={nodes near coords, point meta=explicit symbolic, visualization depends on={value \thisrow{anchor}\as\lblanchor},
    every node near coord/.append style={font=\sffamily\tiny, anchor=\lblanchor, text=ink2, inner sep=1.5pt}},
]
\fill[sone!7] (axis cs:3,1) rectangle (axis cs:8,60);
\pgfplotsinvokeforeach{0.1,1,10}{
  \draw[gen!70, dashed, line width=0.4pt] (axis cs:3,#1) -- (axis cs:100,#1);
  \node[font=\sffamily\tiny, text=gen!80!black, anchor=north east, inner sep=1pt] at (axis cs:100,#1) {#1\,Gbit/s link};}
\addplot[only marks, mark=*, mark size=0.8pt, gen!40, discard if not={group}{general}] table[x=bits, y=gbit] {\gen{fignet_roof_pts.tsv}};
\addplot[only marks, mark=*, mark size=0.8pt, dom!40, discard if not={group}{domain}] table[x=bits, y=gbit] {\gen{fignet_roof_pts.tsv}};
\addplot[only marks, mark=*, mark size=0.8pt, sone!45, discard if not={group}{s1}] table[x=bits, y=gbit] {\gen{fignet_roof_pts.tsv}};
\addplot[only marks, lbl, mark=*, mark size=1.9pt, gen!80!black, discard if not={whole}{0}, discard if not y={group}{general}]
  table[x=bits, y=gbit, meta=label] {\gen{fignet_roof.tsv}};
\addplot[only marks, lbl, mark=square*, mark size=1.8pt, gen!80!black, discard if not={whole}{1}]
  table[x=bits, y=gbit, meta=label] {\gen{fignet_roof.tsv}};
\addplot[only marks, lbl, mark=*, mark size=1.9pt, dom, discard if not={group}{domain}] table[x=bits, y=gbit, meta=label] {\gen{fignet_roof.tsv}};
\addplot[only marks, lbl, mark=*, mark size=2.8pt, sone, discard if not={group}{s1}] table[x=bits, y=gbit, meta=label] {\gen{fignet_roof.tsv}};
\node[font=\sffamily\tiny, text=sone, anchor=north west, align=left] at (axis cs:3.1,55) {small and keeps up:\\the network sets the rate};
\end{axis}
\end{tikzpicture}

\vspace{4pt}
% ---------------- (c) measured readers against the two ceilings ----------------
\begin{tikzpicture}
\begin{axis}[
  title={c\quad Measured readers, \nm{fig1b-dataset}},
  width=0.47\linewidth, height=6.4cm, xmode=log, ymode=log, xmin=0.03, xmax=30, ymin=50, ymax=400000,
  xlabel={link bandwidth (Gbit/s), 1\,ms latency}, ylabel={cells per second into the model (sustained)},
  xtick={0.1,1,10}, xticklabels={0.1,1,10}, axis x line*=bottom, axis y line*=left,
  legend style={font=\sffamily\tiny, draw=none, fill=none, at={(0.5,-0.24)}, anchor=north, legend columns=2, cells={anchor=west}},
]
\addplot[sone, dashed, no marks, forget plot] table[x=gbit, y=net] {\gen{fignet_ceilings.tsv}}
  node[pos=0.3, sloped, above, font=\sffamily\tiny, text=sone] {network ceiling for \sone};
\addplot[ink2, dashed, no marks, forget plot] table[x=gbit, y=dec] {\gen{fignet_ceilings.tsv}}
  node[pos=0.02, anchor=south west, font=\sffamily\tiny, text=ink2] {decoder ceiling, \nm{threadsHeadline} cores};
\addplot[sone, very thick, mark=*, mark size=1.8pt, discard if not={loader}{s1_stream}] table[x=gbit, y=cells_per_s] {\gen{fignet_loaders.tsv}};
\addplot[sone!55, thick, dotted, mark=o, mark size=1.6pt, discard if not={loader}{s1_stream_seq}] table[x=gbit, y=cells_per_s] {\gen{fignet_loaders.tsv}};
\addplot[dec, mark=triangle*, mark size=1.6pt, discard if not={loader}{anndata_zarr_lazy}] table[x=gbit, y=cells_per_s] {\gen{fignet_loaders.tsv}};
\addplot[dom, mark=diamond*, mark size=1.6pt, discard if not={loader}{soma_ml}] table[x=gbit, y=cells_per_s] {\gen{fignet_loaders.tsv}};
\legend{\textbf{.s1}, {.s1, one request at a time}, AnnData Zarr, TileDB-SOMA ML}
\end{axis}
\end{tikzpicture}\hfill
% ---------------- (d) first minibatch against latency ----------------
\begin{tikzpicture}
\begin{axis}[
  title={d\quad Round trips set the first minibatch},
  width=0.47\linewidth, height=6.4cm, xmin=0, xmax=100, ymin=0,
  xlabel={round-trip latency (ms), 1\,Gbit/s link}, ylabel={seconds to the first minibatch},
  axis x line*=bottom, axis y line*=left,
  legend style={font=\sffamily\tiny, draw=none, fill=none, at={(0.02,0.98)}, anchor=north west, cells={anchor=west}},
]
\fill[gen!10] (axis cs:10,0) rectangle (axis cs:100,1.8);
\addplot[sone, thick, no marks, forget plot, discard if not={loader}{s1_stream}, discard if not y={measured}{0}] table[x=rtt_ms, y=first_s] {\gen{fignet_latency.tsv}};
\addplot[sone!55, thick, dotted, no marks, forget plot, discard if not={loader}{s1_stream_seq}, discard if not y={measured}{0}] table[x=rtt_ms, y=first_s] {\gen{fignet_latency.tsv}};
\addplot[dec, no marks, forget plot, discard if not={loader}{anndata_zarr_lazy}, discard if not y={measured}{0}] table[x=rtt_ms, y=first_s] {\gen{fignet_latency.tsv}};
\addplot[sone, only marks, mark=*, mark size=2pt, discard if not={loader}{s1_stream}, discard if not y={measured}{1}] table[x=rtt_ms, y=first_s] {\gen{fignet_latency.tsv}};
\addplot[sone!55, only marks, mark=o, mark size=2pt, discard if not={loader}{s1_stream_seq}, discard if not y={measured}{1}] table[x=rtt_ms, y=first_s] {\gen{fignet_latency.tsv}};
\addplot[dec, only marks, mark=triangle*, mark size=2pt, discard if not={loader}{anndata_zarr_lazy}, discard if not y={measured}{1}] table[x=rtt_ms, y=first_s] {\gen{fignet_latency.tsv}};
\legend{{\textbf{.s1}: \nm{net-rtts-s1_stream} round trips in sequence}, {.s1, one request at a time: \nm{net-rtts-s1_stream_seq}}, {AnnData Zarr: \nm{net-rtts-anndata_zarr_lazy}}}
\node[font=\sffamily\tiny, text=gen!80!black, anchor=south west, align=left] at (axis cs:12,0.02) {cloud object stores:\\tens of ms to a first byte};
\end{axis}
\end{tikzpicture}
\caption{\textbf{What reading at the speed of the network means, and which formats and readers do.}
\textbf{(a)} Non-zeros delivered per second over a link of bandwidth $B$, $R_{\mathrm{eff}}=(b/B+1/D_p)^{-1}$
(Eq.~\ref{eq:reff}), for every format, from its median size $b$ and in-memory read rate $D_p$ on \nm{threadsHeadline}
cores over the matrices. On the left of its knee a format is network-bound and delivers $B/b$ (dotted diagonal for
\sone): the fewer bits per non-zero, the more non-zeros the same link carries. On the right it is decoder-bound at $D_p$
(dotted horizontal). Grey vertical lines: the tiers measured in panel c.
\textbf{(b)} The same formats by size and by the stored bits their reader consumes per second, $D_p b$: a format keeps the
network as the limit on every link below its point. Large points, medians; small points, single matrices; squares,
formats that cannot be read in part. Shaded: under 8 bits per non-zero and fast enough for a 1\,Gbit/s link.
\textbf{(c)} Sustained cells per second delivered to the model from \nm{fig1b-dataset} through the object store, against the
two ceilings for \sone: the link's bandwidth divided by the file's \nm{net-ds-bytes-per-cell} bytes per cell (dashed
diagonal), and the decoder's in-memory rate on \nm{threadsHeadline} cores (dashed horizontal).
\textbf{(d)} Seconds to the first 512-cell minibatch at 1 and 30\,ms round-trip latency (points; median over the four
panel datasets) and the straight line through them, whose slope counts the round trips a reader pays one after another.}
\label{fig:net}
\end{figure*}


\subsection{\sone{} is close to the limits of size and speed}
\label{sec:floor}

How much smaller could a file be? Following the payload of a matrix as ideas are added (Fig.~\ref{fig:floor}a), storing
the steps between entries and compressing them with the strongest general-purpose compressor shrinks the file a lot, but
modelling the counts, as \sone{} does, removes nearly a third of what is left. The best count model, even though it was
charged nothing for its own parameters, was only a few percent shorter than \sone. For data counted with molecular
barcodes, most of the remaining bits describe which molecules happened to be captured, a random sample that no lossless
format can avoid storing. Nor is anything left for a byte-level compressor: zstd-19 does not shrink an \sone{} bitstream
at all, and a code that treats each step and value on its own cannot go below \nm{ent-h0-vs-floor} times the floor. In
size, any lossless format can improve on \sone{} by single-digit percentages at most.

How fast could a file be decoded? Timed on one core, the integer codes trade size for speed along a clear frontier
(Fig.~\ref{fig:floor}b). Byte- and bit-aligned codes decode at close to the rate at which a core can merely write the
output, but store three to twelve times more bits than \sone; entropy coders reach small sizes at a fraction of that
speed. \sone{} is the smallest end of this frontier, and no code we measured was both as small and faster. The ideal code,
as small as the best count model and as fast as writing the output, lies far from every measured code.

Is \sone{} slow for its size, or is that the price of being small? The time the other codes spend per entry, beyond
writing the output, grows as a power of how few bits they store. Extrapolated to \sone's size, this trend predicts the rate
a code that small could be expected to reach, and \sone{} decodes at \nm{trend-s1-pct}\% of it (95\% interval
\nm{trend-s1-lo}--\nm{trend-s1-hi}\%): within the uncertainty, it lies on the trend. Reading genes in their source order
instead of the stored order would put it at \nm{trend-source-pct}\%.

What matters in practice is the rate a link delivers (Eq.~\ref{eq:reff}), and relative to the ideal code it depends on how
much bandwidth each decoding core must keep up with (Fig.~\ref{fig:floor}c). When cores are plentiful relative to the
link, as when a 32-core machine reads from a 1\,Gbit/s connection, \sone{} delivers \nm{frac-s1-32c}\% of what the ideal
code would and more than any other code, because size is what limits. When few cores serve a fast link, the decoder
becomes the limit: above about \nm{frac-cross-gbit}\,Gbit/s per core, zstd's larger but faster stream delivers more.
Decoding speed, not size, is where the remaining headroom lies.

% Figure 4: how close .s1 is to the limits of size and decoding speed, and what it costs to write. Values from
% generated/fig3a_ladder.tsv (ladder stage) and generated/figlim_*.tsv (kernels and writecost stages, report_limits.py).
% More: fig_supp_floor.tex.
\begin{figure*}[t]
\centering
\pgfmathsetlengthmacro{\flX}{3.35cm}            % left of a and c: a's row labels
\pgfmathsetlengthmacro{\flA}{5.7cm}             % side of every (square) panel
\pgfmathsetlengthmacro{\flG}{1.9cm}             % between the columns: the y labels of b and d
\pgfmathsetlengthmacro{\flB}{\flA}
\pgfmathsetlengthmacro{\flH}{\flA}
\pgfmathsetlengthmacro{\flHH}{\flA}
\pgfmathsetlengthmacro{\flV}{1.65cm}            % between rows
\begin{tikzpicture}
\pgfplotsset{every axis title/.append style={at={(0,1)}, anchor=south west, xshift=-4pt}}
% ---------------- (a) bits per non-zero as each idea is added ----------------
\begin{axis}[name=fa, at={(\flX,0)}, anchor=north west, scale only axis, width=\flA, height=\flH,
  title={a\quad From CSC to \sone, payload only},
  xbar, bar shift=0pt, bar width=0.5,
  xmin=0, xlabel={bits per non-zero (row indices and values), median},
  ytick=data, ytick style={draw=none}, yticklabels from table={\gen{fig3a_ladder.tsv}}{label}, y dir=reverse,
  ymin=-0.5, ymax=5.5,
  axis x line*=bottom, axis y line*=left,
  nodes near coords, every node near coord/.append style={font=\sffamily\tiny, anchor=west, /pgf/number format/fixed, /pgf/number format/precision=1},
  error bars/error bar style={ink2, line width=0.5pt}, error bars/error mark options={rotate=90, mark size=1.5pt, ink2},
]
\addplot[draw=none, forget plot, nodes near coords={}] table[x=median, y=pos] {\gen{fig3a_ladder.tsv}};
\addplot[fill=gen!70, draw=none, discard if not={key}{L0_csc_i32}] table[x=median, y=pos] {\gen{fig3a_ladder.tsv}};
\addplot[fill=gen!70, draw=none, discard if not={key}{L1_leb128}] table[x=median, y=pos] {\gen{fig3a_ladder.tsv}};
\addplot[fill=gen!70, draw=none, discard if not={key}{L2_gaps_leb_zstd3}] table[x=median, y=pos] {\gen{fig3a_ladder.tsv}};
\addplot[fill=gen!70, draw=none, discard if not={key}{L3_gaps_leb_zstd19}] table[x=median, y=pos] {\gen{fig3a_ladder.tsv}};
\addplot[fill=sone, draw=none, discard if not={key}{L4_s1_payload}, error bars/.cd, x dir=both, x explicit]
  table[x=median, y=pos, x error plus=ep, x error minus=em] {\gen{fig3a_ladder.tsv}};
\addplot[fill=white, draw=ink2, dashed, discard if not={key}{floor}, error bars/.cd, x dir=both, x explicit]
  table[x=median, y=pos, x error plus=ep, x error minus=em] {\gen{fig3a_ladder.tsv}};
\end{axis}
% ---------------- (b) size against decoding speed on one core ----------------
\begin{axis}[name=fb, at={($(fa.north east)+(\flG,0)$)}, anchor=north west, scale only axis, width=\flB, height=\flH,
  title={b\quad Size against decoding speed, one core},
  xmode=log, ymode=log, xmin=3, xmax=80, ymin=20, ymax=2500,
  xtick={4,5,10,20,50}, xticklabels={4,5,10,20,50}, ytick={20,50,100,200,500,1000,2000}, yticklabels={20,50,100,200,500,1000,2000},
  xlabel={stored bits per non-zero}, ylabel={decoded per core, million non-zeros/s},
  axis x line*=bottom, axis y line*=left, clip mode=individual,
  lbl/.style={nodes near coords, point meta=explicit symbolic, visualization depends on={value \thisrow{anchor}\as\lblanchor},
    every node near coord/.append style={font=\sffamily\tiny, anchor=\lblanchor, text=ink2, inner sep=2pt}},
]
% lines of equal input rate: one core decoding on such a line keeps a link of that speed full (b D = c Gbit/s)
\pgfplotsinvokeforeach{0.3,1,3,10}{
  \addplot[domain=3:{min(80,#1*1000/21)}, samples=2, gen!45, densely dotted, forget plot] {#1*1000/x}
    node[pos=0.85, anchor=south, font=\sffamily\tiny, text=gen, sloped, inner sep=1pt] {#1\,Gbit/s};
}
\addplot[only marks, mark=*, mark size=0.6pt, gen!80!black, opacity=0.35, forget plot, discard if not={group}{other}] table[x=bits, y=mnnzs] {\gen{figlim_plane_pts.tsv}};
\addplot[only marks, mark=*, mark size=0.7pt, sone, opacity=0.5, forget plot, discard if not={group}{s1}] table[x=bits, y=mnnzs] {\gen{figlim_plane_pts.tsv}};
% the trend of the other codes, fitted on their frontier and extrapolated to the floor
\addplot[tens, dashed, line width=0.8pt, no marks, forget plot] table[x=bits, y=mnnzs] {\gen{figlim_trend.tsv}};
\node[anchor=north west, font=\sffamily\tiny, text=tens, inner sep=1pt, align=left] at (axis cs:11.5,300) {trend of the\\other codes};
\addplot[only marks, mark=diamond, mark size=2.4pt, tens, thick, fill=white, forget plot] table[x=bits, y=mnnzs] {\gen{figlim_trend_floor.tsv}};
% the floor (best count model) and the output bound (writing the decoded arrays), and the ideal code where they meet
\addplot[dl!80!black, dashed, line width=0.6pt, no marks, forget plot] table[x=bits, y=mnnzs] {\gen{figlim_floorline.tsv}}
  node[pos=0.03, anchor=south west, rotate=90, font=\sffamily\tiny, text=dl!80!black] {best count model};
\addplot[dec!80!black, dashed, line width=0.6pt, no marks, forget plot] table[x=bits, y=mnnzs] {\gen{figlim_storeline.tsv}}
  node[pos=0.1, anchor=north west, font=\sffamily\tiny, text=dec!80!black] {writing the output only};
\addplot[only marks, mark=star, mark size=3.5pt, ink2, thick, forget plot] table[x=bits, y=mnnzs] {\gen{figlim_ideal.tsv}}
  node[anchor=south west, font=\sffamily\tiny, text=ink2, xshift=1pt] {ideal: model size, output speed};
\addplot[only marks, lbl, mark=*, mark size=1.9pt, gen, draw=white, line width=0.3pt, discard if not={group}{fast}]
  table[x=bits, y=mnnzs, meta=label] {\gen{figlim_plane.tsv}};
\addplot[only marks, lbl, mark=*, mark size=1.9pt, dec, draw=white, line width=0.3pt, discard if not={group}{entropy}]
  table[x=bits, y=mnnzs, meta=label] {\gen{figlim_plane.tsv}};
\addplot[only marks, lbl, mark=*, mark size=2.6pt, sone, draw=white, line width=0.3pt, discard if not={group}{s1}]
  table[x=bits, y=mnnzs, meta=label] {\gen{figlim_plane.tsv}};
\end{axis}
% ---------------- (c) share of the ideal code's delivered rate, against bandwidth per decoding core ----------------
\begin{axis}[name=fc, at={($(fa.south west)+(0,-\flV)$)}, anchor=north west, scale only axis, width=\flA, height=\flHH,
  title={c\quad Delivered rate, share of the ideal code},
  xmode=log, xmin=0.01, xmax=10, ymin=0, ymax=1.02,
  xtick={0.01,0.1,1,10}, xticklabels={0.01,0.1,1,10}, ytick={0,0.25,0.5,0.75,1}, yticklabels={0,25,50,75,100\%},
  xlabel={link bandwidth per decoding core (Gbit/s)}, ylabel={share of the ideal code's rate},
  axis x line*=bottom, axis y line*=left, clip=false,
  legend style={font=\sffamily\tiny, draw=none, fill=white, fill opacity=0.85, text opacity=1, at={(0.02,0.03)},
    anchor=south west, legend columns=1, cells={anchor=west}},
]
\addplot[sone, line width=1.6pt, no marks] table[x=x, y=pz2_avx2] {\gen{figlim_frac.tsv}};
\addplot[sone, dashed, thick, no marks] table[x=x, y=practical] {\gen{figlim_frac.tsv}};
\addplot[dec, thick, no marks] table[x=x, y=zstd19_leb_gaps] {\gen{figlim_frac.tsv}};
\addplot[ycomb, gen!60, line width=0.4pt, no marks, forget plot] table[x=x, y expr=1.02] {\gen{figlim_marks.tsv}};
\addplot[only marks, mark=none, nodes near coords, point meta=explicit symbolic, forget plot,
  every node near coord/.append style={font=\sffamily\tiny, text=gen, rotate=90, anchor=north east, inner sep=1.5pt}]
  table[x=x, y expr=1.0, meta=label] {\gen{figlim_marks.tsv}};
\legend{\sone, {\sone's size at zstd's speed}, {zstd-19}}
\end{axis}
% ---------------- (d) the cost of writing, as fold changes over writing uncompressed CSC ----------------
\begin{axis}[name=fd, at={($(fb.south west)+(0,-\flV)$)}, anchor=north west, scale only axis, width=\flB, height=\flHH,
  title={d\quad The cost of writing, relative to plain CSC arrays},
  xmode=log, ymode=log, xmin=0.5, xmax=2000, ymin=0.7, ymax=40,
  xtick={1,10,100,1000}, xticklabels={1$\times$,10$\times$,100$\times$,1000$\times$},
  ytick={1,3,10,30}, yticklabels={1$\times$,3$\times$,10$\times$,30$\times$},
  xlabel={time to write (log scale)}, ylabel={memory beyond the input},
  axis x line*=bottom, axis y line*=left, clip=false,
  lbl/.style={nodes near coords, point meta=explicit symbolic, visualization depends on={value \thisrow{anchor}\as\lblanchor},
    every node near coord/.append style={font=\sffamily\tiny, anchor=\lblanchor, text=ink2, inner sep=3.5pt}},
  rng/.style={error bars/x dir=both, error bars/y dir=both, error bars/x explicit, error bars/y explicit,
    error bars/error bar style={line width=0.5pt, opacity=0.55}, error bars/error mark=none},
]
\expandafter\pgfplotsinvokeforeach\expandafter{\otherformats}{
  \addplot[only marks, lbl, mark=*, mark size=1.8pt, c-#1, rng, discard if not={format}{#1}]
    table[x=t, y=m, meta=label, x error plus expr=\thisrow{t_hi}-\thisrow{t}, x error minus expr=\thisrow{t}-\thisrow{t_lo},
      y error plus expr=\thisrow{m_hi}-\thisrow{m}, y error minus expr=\thisrow{m}-\thisrow{m_lo}] {\gen{figlim_write.tsv}};
}
\addplot[only marks, lbl, mark=*, mark size=2.6pt, c-s1, draw=white, line width=0.3pt, rng, discard if not={format}{s1}]
  table[x=t, y=m, meta=label, x error plus expr=\thisrow{t_hi}-\thisrow{t}, x error minus expr=\thisrow{t}-\thisrow{t_lo},
    y error plus expr=\thisrow{m_hi}-\thisrow{m}, y error minus expr=\thisrow{m}-\thisrow{m_lo}] {\gen{figlim_write.tsv}};
\end{axis}
\end{tikzpicture}
\caption{\textbf{How close \sone{} is to the limits of size and speed, and what it costs.}
\textbf{(a)} Payload (row indices and values, without names) in bits per non-zero as each idea is added, median over
matrices; whiskers, interquartile range. Dashed: the floor, the shortest code length of four count models fitted on half of
each matrix's cells and scored on the other half (a reference, not a bound); it is a median of \nm{s1AboveFloorPct}\%
below \sone.
\textbf{(b)} Integer codes (grey: byte- and bit-aligned; yellow: entropy-coded) and \sone, each encoding the matrix's own
row indices and values and decoding them on one core, checked bit for bit; \nm{kern-n} matrices, faint points single
matrices, large points medians. Dotted: the link speed one core keeps up with. Green: the floor; yellow: the rate of
only writing the decoded arrays. Star: the ideal code, where the two meet. Purple dashed: the trend of the fastest other
codes at each size, extrapolated to the floor (diamond); \sone{}
decodes at \nm{trend-s1-pct}\% of the rate it predicts for \sone's size (95\% interval
\nm{trend-s1-lo}--\nm{trend-s1-hi}\%).
\textbf{(c)} Non-zeros per second a link delivers (Eq.~\ref{eq:reff}, decoder on all cores) relative to the ideal code of
(b), against the link bandwidth each decoding core serves; median of per-matrix ratios. Dashed: a code as small as
\sone{} that decoded as fast as zstd. Grey ticks mark familiar set-ups.
\textbf{(d)} Time to write a matrix held in memory, and the memory each writer needs beyond its input, each relative to
writing uncompressed CSC arrays of the same matrix; median over \nm{figlim-write-n} matrices of
\nm{figlim-write-cells-lo} to \nm{figlim-write-cells-hi} cells, lines their range.}
\label{fig:floor}
\end{figure*}


\subsection{Costs and use}

Writing is the price of a small file (Fig.~\ref{fig:floor}d). \sone{} wrote a 580,000-cell atlas faster than
gzip-compressed AnnData, the 10x HDF5 file or TileDB-SOMA, but slower than formats that store raw or lightly compressed
arrays, such as Zarr, BPCells and Parquet, and its writer needed about three times as much memory beyond the input as
AnnData's.

\sone{} files are written and read by the open-source \texttt{singlet} package (Fig.~\ref{fig:eco}). One call stores a
matrix from AnnData, SciPy, SingleCellExperiment or Seurat, and one call loads it back as any of these, with exactly the
original counts, gene identifiers and barcodes; loading was faster than reading the equivalent AnnData or 10x file in both
Python and R. For training, \texttt{singlet.torch.S1Cells} is a PyTorch dataset that streams shuffled minibatches from a
local path or a URL.

% Figure 4: .s1 in the ecosystem. Flow chart with the calls a user runs (all in singlet 1.3.x).
\begin{figure*}[tbp]
\centering
\resizebox{\linewidth}{!}{%
\begin{tikzpicture}[
  font=\sffamily\scriptsize,
  src/.style={draw=gen!80!black, fill=gen!10, rounded corners=1.5pt, minimum width=2.9cm, minimum height=0.62cm, align=center},
  hub/.style={draw=sone!80!black, fill=sone!12, rounded corners=2pt, align=center, inner sep=4pt},
  out/.style={draw=#1!80!black, fill=#1!10, rounded corners=1.5pt, minimum width=3.3cm, minimum height=0.62cm, align=center},
  code/.style={draw=ink2!40, fill=ink2!5, rounded corners=1pt, align=left, inner sep=3pt, font=\ttfamily\tiny, anchor=west},
  arr/.style={-{Stealth[length=4pt]}, thick, ink2},
  hdr/.style={font=\sffamily\small\bfseries, text=ink2},
]
% ---------- column headers ----------
\node[hdr] at (0,0.75) {Your data};
\node[hdr] at (5.6,0.75) {\sone{} files};
\node[hdr] at (11.0,0.75) {Where you work};

% ---------- sources ----------
\node[src] (cr) at (0,0)     {Cell Ranger HDF5 / MTX};
\node[src] (ss) at (0,-0.85) {STARsolo output};
\node[src] (ad) at (0,-1.7)  {AnnData (.h5ad, Zarr)};
\node[src] (sr) at (0,-2.55) {SCE / Seurat (R)};
\node[src] (sp) at (0,-3.4)  {spatial count matrices};

% ---------- writer ----------
\node[hub, minimum width=2.6cm] (wr) at (2.95,-1.7) {\textbf{singlet writer}\\lossless, verified\\on write};
\foreach \s in {cr,ss,ad,sr,sp} \draw[arr] (\s.east) -- (wr.west);

% ---------- the objects ----------
\node[hub, minimum width=2.9cm, minimum height=2.3cm] (obj) at (5.6,-1.7) {};
\node[font=\sffamily\scriptsize\bfseries, text=sone!80!black] at (5.6,-0.82) {disk or object store};
\node[draw=sone, fill=white, minimum width=2.4cm, font=\sffamily\tiny] at (5.6,-1.2) {dataset manifest (.json)};
\node[draw=sone, fill=sone!30, minimum width=2.4cm, font=\sffamily\tiny] at (5.6,-1.6) {sample objects (.s1)};
\node[draw=gen, fill=gen!15, minimum width=2.4cm, font=\sffamily\tiny] at (5.6,-2.0) {packs: genes, barcodes};
\node[font=\sffamily\tiny, text=ink2] at (5.6,-2.5) {HTTP range reads};
\draw[arr] (wr.east) -- (obj.west);

% ---------- readers ----------
\node[out=sone] (py)  at (10.0,0.0)  {\textbf{AnnData} (Python)};
\node[out=tens] (pt)  at (10.0,-1.0) {\textbf{PyTorch} sparse minibatches};
\node[out=dl]   (sce) at (10.0,-2.0) {\textbf{SingleCellExperiment} (R)};
\node[out=dl]   (seu) at (10.0,-3.0) {\textbf{Seurat} (R)};
\node[out=gen]  (cpp) at (10.0,-4.0) {\textbf{C++} CSC arrays};
\foreach \o in {py,pt,sce,seu,cpp} \draw[arr] (obj.east) -- (\o.west);

% ---------- code beside each reader ----------
\node[code] at (11.8,0.0)  {import singlet\\adata = singlet.load("pbmc10k.s1")};
\node[code] at (11.8,-1.0) {from singlet.torch import S1Cells\\ds = S1Cells("https://.../pbmc10k.s1", batch\_size=512)\\for b in DataLoader(ds, batch\_size=None): b["X"]};
\node[code] at (11.8,-2.0) {sce <- singlet::load("pbmc10k.s1", as = "sce")};
\node[code] at (11.8,-3.0) {so \ <- singlet::load("pbmc10k.s1", as = "seurat")};
\node[code] at (11.8,-4.0) {auto m = pz2::Matrix::open(src);\\m.decode\_columns(src, c0, c1, out);};

% ---------- writer calls ----------
\node[code] at (-1.45,-4.6) {singlet.write\_s1(adata, "sample.s1")\ \ \ \ \ \ \ \ \# Python\\singlet::write\_s1(sce, "sample.s1")\ \ \ \ \ \ \ \ \ \# R};
\end{tikzpicture}}
\caption{\textbf{\sone{} in the single-cell ecosystem.} Count matrices from Cell Ranger, STARsolo, AnnData,
SingleCellExperiment, Seurat or spatial pipelines are written once as \sone{} files, and the writer checks that the file
decodes to exactly the input. The \texttt{singlet} package reads them into AnnData, PyTorch minibatches,
SingleCellExperiment, Seurat or plain C++ arrays, from a local file or by HTTP range requests. The code beside each
target is the complete call (\texttt{DataLoader} is PyTorch's); Table~\ref{tab:eco} gives the time each takes on two
benchmark datasets. All calls are in \texttt{singlet}~\nm{rel-singlet_py} (Python) and \nm{rel-singlet_r} (R).}
\label{fig:eco}
\end{figure*}

